% Figure: temperature dependence of the semiconductor band gap (Varshni law) for
% silicon and gallium arsenide. The gap shrinks as the lattice expands and the
% electron-phonon coupling grows. Sample points from the Varshni relation
% E_g(T) = E_g(0) - aT^2/(T+b) computed offline and entered explicitly.
\begin{figure}[htbp]
\centering
\begin{tikzpicture}[font=\scriptsize]
\begin{axis}[
    width=10cm, height=6cm,
    xlabel={temperature $T$ (\si{\kelvin})},
    ylabel={band gap $E_g$ (\si{\electronvolt})},
    xmin=0, xmax=600,
    ymin=1.0, ymax=1.6,
    xtick={0,100,200,300,400,500,600},
    ytick={1.0,1.1,1.2,1.3,1.4,1.5,1.6},
    grid=both,
    grid style={enggray!25},
    tick label style={font=\tiny},
    label style={font=\scriptsize},
    axis line style={enggray},
    legend style={font=\tiny, at={(0.97,0.97)}, anchor=north east,
                  draw=enggray!50},
    every axis plot/.append style={very thick},
]
% GaAs: Eg(0)=1.519, a=5.41e-4, b=204
\addplot[engblue, mark=*, mark size=1.0pt] coordinates {
  (0,1.519) (50,1.517) (100,1.510) (150,1.501) (200,1.489)
  (250,1.476) (300,1.462) (350,1.447) (400,1.431) (450,1.415)
  (500,1.398) (550,1.381) (600,1.363)
};
\addlegendentry{GaAs}
% Si: Eg(0)=1.170, a=4.73e-4, b=636
\addplot[engred, dashed, mark=square, mark size=1.0pt] coordinates {
  (0,1.170) (50,1.169) (100,1.163) (150,1.156) (200,1.147)
  (250,1.137) (300,1.125) (350,1.114) (400,1.101) (450,1.089)
  (500,1.076) (550,1.063) (600,1.050)
};
\addlegendentry{Si}
% mark room-temperature values
\addplot[engorange, only marks, mark=o, mark size=2.6pt]
  coordinates {(300,1.462) (300,1.125)};
\node[engorange, anchor=south, font=\tiny] at (axis cs:300,1.475)
  {$300\,\si{\kelvin}$};
\end{axis}
\end{tikzpicture}
\caption[Band gap against temperature]{Band gap of silicon and gallium arsenide
against temperature, the Varshni relation $E_g(T)=E_g(0)-aT^{2}/(T+b)$. The gap
shrinks as temperature rises because the lattice expands and the electron-phonon
interaction strengthens, a drop of roughly $0.3$ to $0.4\,\si{\milli
\electronvolt}$ per kelvin near room temperature. Two device consequences ride
on this slope. The intrinsic carrier density climbs steeply because its
exponent carries $E_g/2k_BT$ with a falling numerator over a rising denominator,
which sets the high-temperature limit of silicon electronics. The emission
wavelength of an LED or laser diode drifts to the red as the device heats, the
wavelength shift the laser-diode case study charged to temperature and the
reason a wavelength-critical source is held at a set temperature.}
\label{fig:vol04ch12:varshni-gap}
\end{figure}
